% ============================================================
%  Chapter: Emotional Intelligence and Attitude
%  Foundational Values recap, Emotional Intelligence, and Attitude
% ============================================================

\chapter{Emotional Intelligence and Attitude}


\vspace{14pt}

\section{Recap: Foundational Values for Civil Services}

The previous class dealt with the foundational values for civil services. This session opens by revising those ideas before moving into fresh examples and new topics.

\subsection{Dedication to Public Service (Passionate Commitment)}

Dedication to public service is very important, and it is essentially a passionate commitment. Passionate commitment means going beyond doing the minimum. That is what dedication to public service means.

\begin{KeyConcept}
\begin{itemize}
\item Passionate commitment = going beyond doing the minimum required of the role.
\item This is the essence of ``dedication to public service''.
\end{itemize}
\end{KeyConcept}

\subsection{Sympathy, Empathy and Compassion}

These are three different things, and they were distinguished carefully:

\begin{itemize}
\item Sympathy: operates mostly at the cognitive level --- you merely understand that the other person is in difficulty. You feel some pity or ``daya'' (mercy) towards them.
\item Empathy: you actually feel the pain of others --- you experience their pain yourself.
\item Compassion: this is when empathy is accompanied by a desire to alleviate the problems of others. There is a desire to make an effort to reduce the other person's pain.
\end{itemize}

Of course, compassion is the biggest of the three, because it carries the desire to actually try to reduce the other person's pain. That is what administration should look for.

\begin{ComparisonBox}
\begin{itemize}
\item Sympathy $\rightarrow$ cognitive understanding + pity (``daya'').
\item Empathy $\rightarrow$ you feel the other's pain.
\item Compassion $\rightarrow$ empathy + desire to elevate/alleviate the other's problems (the highest).
\end{itemize}
\end{ComparisonBox}

Various examples were discussed for how administration should be empathetic and compassionate.

\paragraph{Weberian Bureaucracy --- the contrast} This was discussed especially in the context of Weberian bureaucracy. Weber asked bureaucracy to be objective only and rational only. Weberian bureaucracy was supposed to be value-neutral --- that is, taking decisions on the basis of facts and facts alone, and not focusing on emotions in decision-making.

\begin{KeyConcept}
\begin{itemize}
\item Weberian bureaucracy = value-neutral: decisions on facts alone, not emotions.
\item It was meant to be strictly objective and rational.
\end{itemize}
\end{KeyConcept}

\subsection{Non-partisanship and Impartiality}

\paragraph{Non-partisanship} Non-partisanship means that we will not align with any political party or ideology. It is important for bureaucrats not to align with a political party or ideologies, because bureaucrats are the permanent executive. On the other hand, political executives are temporary in nature --- they keep coming and going, because governments can change.

Governments can change even during the tenure in which these bureaucrats were recruited. When you become an IAS officer, you will serve the Government of India for roughly 25--35 years, and it is not certain which government will remain in power over that period.

It is important that every new government formed at the central level should be able to trust the advice of the civil servant, because civil servants come with expertise and a lot of experience. They should give this advice both in policy-making and in policy implementation. When advice is given to ministers, ministers should not feel that the officer is a ``Congress person'' or a ``BJP person''. They should believe in the objectivity of the advice, and should also trust the civil servant with the implementation of their policies. For example, in implementing ``Make in India'', it is bureaucrats who are the first implementers. That is why it becomes very important that administrators can be trusted. For policy formulation, for objective advice, and for implementation of framed policies, non-partisanship is essential.

\begin{CriticalPoint}
\begin{itemize}
\item Non-partisanship is a NEGATIVE concept: it bars any association with any political party or ideology.
\item Rationale: bureaucrats are the permanent executive; political executives are temporary.
\item Needed for (a) objective policy advice, and (b) faithful implementation across changing governments.
\end{itemize}
\end{CriticalPoint}

If transfers, postings and promotions begin to happen on a partisan basis --- on the basis of alignment with a political party --- then who will work hard? What will be the motivation and morale of the civil services? Additionally, citizens will lose faith in the administration if there is no non-partisanship, because they will feel the officer is ``a BJP person'' or ``a Congress person'' or ``a party person''. When you start behaving as a stooge of a political party rather than as a bureaucrat following constitutional morality, people lose faith in the administration.

\begin{ExampleBox}
Land acquisition and conflict resolution were given as cases where non-partisanship becomes very important, so that the administrator can be trusted by the people.
\end{ExampleBox}

\paragraph{Impartiality} Impartiality means that whenever you are engaging with anyone, you should not discriminate amongst them --- for example, on the basis of race, caste, sex, religion, etc. Impartiality is a positive concept, because it is not barring your association with anyone; it is in fact asking that whenever you interact with anyone, you should not discriminate amongst them.

\begin{ComparisonBox}
\begin{itemize}
\item Non-partisanship = negative concept: bars association with any political party.
\item Impartiality = positive concept: does not bar association with anyone; instead demands non-discrimination during any interaction.
\end{itemize}
\end{ComparisonBox}

A simple example: as the Election Commission of India, if I have two states or two political parties before me and I am not impartial, I would favour one. Impartiality is more of a positive concept --- you remain impartial across political parties, saying in effect, ``I am not here to take sides.''

\begin{QuickRevision}
\begin{itemize}
\item Passionate commitment = beyond the minimum.
\item Sympathy (cognitive/pity) $<$ Empathy (feel the pain) $<$ Compassion (feel + desire to alleviate).
\item Weberian bureaucracy = value-neutral, facts alone.
\item Non-partisanship = negative concept (no party alignment); Impartiality = positive concept (no discrimination).
\end{itemize}
\end{QuickRevision}

\section{Examples \& PYQ Themes on Foundational Values}

New examples and past-year-question (PYQ) themes were introduced, showing where UPSC frames questions.

\subsection{Objectivity without Compassion --- the Mining Districts Case}

Over half of the one lakh crore rupees collected for the welfare of the mining districts was not spent. Funds are often diverted to activities that are not directly linked to the welfare of the mining district. If you do not know what to do, then why collect so much for the welfare of the worker? This is objectivity but without compassion --- and in this context empathy and compassion become much more important.

\begin{DataFact}
\begin{itemize}
\item Over half of \rupee 1 lakh crore collected for welfare of mining districts was not spent.
\item Funds often diverted to activities not directly linked to mining-district welfare.
\item Illustrates objectivity WITHOUT compassion.
\end{itemize}
\end{DataFact}

\subsection{PYQ Theme --- Integrity and Knowledge}

Quotation-based theme: ``Integrity without knowledge is weak and useless, but knowledge without integrity is dangerous and dreadful.'' When we talk about integrity we are really talking about two important axes: one is integrity, and another is knowledge. This is exactly the place where UPSC frames questions.

\begin{PYQLink}
\begin{itemize}
\item PYQ theme: ``Integrity without knowledge is weak and useless, but knowledge without integrity is dangerous and dreadful.''
\item Also: ``In looking for people to hire, you look for three qualities: integrity, intelligence and energy. If they do not have the first, the other two will kill you.''
\item Also: ``Why should impartiality and non-partisanship be considered as foundational values, especially in the socio-political context?''
\end{itemize}
\end{PYQLink}

\paragraph{Integrity without knowledge (weak \& useless)} Consider a surgeon who has integrity but lacks skill, versus a highly skilled surgeon whose success rate is very poor. The example given: a surgeon who is not competent --- around 90\% of patients die and only about 10\% survive the operation. The point is that mere integrity plus incompetence is useless. If doctors are incompetent, if teachers are incompetent, the students suffer. So integrity without knowledge is weak and useless --- for a public servant, competence matters because decisions have consequences for the public.

\paragraph{Knowledge without integrity (dangerous \& dreadful)} On the other dimension, knowledge without integrity is dangerous and dreadful. As noted in the inculcation of values: ``knowledge without values makes a man a more clever devil.'' If a person has knowledge but no integrity/values, that knowledge only makes him a more clever devil --- a sin. This is why the theme repeats: integrity without knowledge is weak and useless, while knowledge without integrity is dreadful and dangerous.

\begin{CriticalPoint}
\begin{itemize}
\item Integrity + no competence $\rightarrow$ weak \& useless (e.g., surgeon with 90\% mortality).
\item Competence + no integrity $\rightarrow$ ``clever devil''; dangerous \& dreadful.
\item Both axes are required in a public servant.
\end{itemize}
\end{CriticalPoint}

\subsection{Impartiality vs.\ Positive Discrimination --- a Caveat}

If impartiality is followed absolutely, it can itself become a source of trouble. Where there is inequality in society, positive discrimination is required --- it should favour the weaker sections. Generalists also make this point: wherever there is inequality in society, positive discrimination should be there to favour weaker sections. So absolutely following impartiality (treating everyone identically regardless of disadvantage) is problematic; some positive discrimination (that is, ``positive partiality'') is justified.

\begin{CriticalPoint}
\begin{itemize}
\item Absolute impartiality can be harmful where society is unequal.
\item Positive discrimination in favour of weaker sections is a justified ``positive partiality''.
\end{itemize}
\end{CriticalPoint}

\section{Ready-to-Use Examples: Apathy, Compassion \& Dedication}

Several current examples were provided that can be used directly in answers.

\paragraph{Examples of Apathy / Lack of Compassion}

\begin{itemize}
\item Sewer deaths: 90\% of the sewer-death workers had no safety gear. If 90\% of deaths involved no safety gear at all, this shows apathy and a lack of compassion. This study was done by the ministry itself --- a social audit commission by the Union Government --- which found that 90\% of sewer deaths involved no safety gear. (Note it briefly.)
\item Forced evictions: People were forcibly evicted from forest land without settlement of their rights under the FRA (Forest Rights Act). The NAPM pointed out this eviction ``without settlement of their rights''. This is an example of apathy --- of how working apathetically goes against social justice.
\end{itemize}

\begin{DataFact}
\begin{itemize}
\item 90\% of sewer-death workers had NO safety gear (social audit commission study by the Union Government).
\item Forest-dwellers forcibly evicted from forest land WITHOUT settlement of their rights under the FRA (raised by NAPM).
\end{itemize}
\end{DataFact}

\paragraph{Examples of Compassion \& Dedication to Public Service}

\begin{itemize}
\item Transgender clinic: India's first clinic for transgender persons reopened in Hyderabad --- a very good example of compassion and dedication, as the first clinic for transgender people.
\end{itemize}

\begin{GovtPolicy}
\begin{itemize}
\item India's first clinic for transgender persons --- reopened in Hyderabad (compassion + dedication).
\item FRA (Forest Rights Act) invoked re: settlement of forest-dwellers' rights before eviction.
\end{itemize}
\end{GovtPolicy}

\begin{QuickRevision}
\begin{itemize}
\item Apathy examples: sewer deaths (90\% no safety gear); forced forest evictions without FRA settlement.
\item Compassion/dedication example: India's first transgender clinic (Hyderabad).
\item Objectivity-without-compassion example: unspent mining-district welfare funds.
\end{itemize}
\end{QuickRevision}

\section{Emotional Intelligence}

\subsection{Opening Quote}

``Anyone can become angry --- that is easy. But to be angry with the right person, to the right degree, at the right time, for the right purpose, and in the right way --- that is not easy.'' The quote conveys that it is easy to be angry; all of us experience emotions. What is important is that emotions come out before the right person, to the right degree, and in the right way. This is what emotional intelligence is.

\begin{KeyConcept}
\begin{itemize}
\item Quote: anger is easy; but anger with the right person, degree, time, purpose and way is hard.
\item That difficulty is exactly what emotional intelligence addresses.
\end{itemize}
\end{KeyConcept}

We all feel emotions, and emotions are a source of energy for us. But how do we regulate those emotions? We should not be driven by emotions. If we are driven by emotion --- for instance if anger keeps rising --- then in anger a person might even commit a hate crime or hurt someone. That is absolutely wrong. You should not be driven by emotions; emotions should not drive you --- you should regulate them. That is emotional intelligence.

\begin{CriticalPoint}
\begin{itemize}
\item Do NOT be driven by emotions (being driven by anger can lead to hate crimes/violence).
\item Emotions should be regulated, not allowed to drive behaviour.
\end{itemize}
\end{CriticalPoint}

\subsection{Definition}

Emotional intelligence is the ability to interpret, understand and manage one's own emotions and the emotions of others.

\paragraph{It is an ``ability'' --- so it can be misused} Look carefully at the definition. First, it is an ability. If it is an ability, it means it can also be misused. Any ability can be misused. For example, a technological ability to read everyone's WhatsApp chats could be used positively (in a security context) or negatively (threatening someone's privacy). The point is that any ability can be misused.

Therefore, if a person is emotionally intelligent, that by itself does not make him a moral person. This links to the ``moral attitude'' idea taught earlier under civil-services attitude: even if you are emotionally intelligent, you will not necessarily use it morally. (The misuse dimension is taken up in detail at the end of the topic.)

\begin{CriticalPoint}
\begin{itemize}
\item Emotional intelligence is an ABILITY $\rightarrow$ it can be misused.
\item Being emotionally intelligent does NOT by itself make a person moral.
\end{itemize}
\end{CriticalPoint}

\paragraph{Interpret/understand + manage = Recognition + Regulation} ``Interpret and understand'' means recognition --- we should be able to recognize, interpret and understand our emotions. ``Manage'' means regulation --- we should be able to control them. Recognition is the first step: if I am getting angry, I should first recognize that I am getting angry; without recognizing it I cannot work on it. Then comes regulation --- controlling it. And this applies not only to our own emotions but also to the emotions of others, because the definition covers ``one's emotions and emotions of others''.

\begin{KeyConcept}
\begin{itemize}
\item Two dimensions: Recognition and Regulation.
\item Recognition = self-awareness (own emotions) + social awareness (others' emotions).
\item Regulation = self-management (own emotions) + relationship management (others' emotions).
\end{itemize}
\end{KeyConcept}

\subsection{Five Components of Emotional Intelligence}

Human beings feel emotions --- anger, frustration, grief --- and much of this is driven by self-motivation. From this flow the five major components of emotional intelligence.

\begin{center}
\begin{tabularx}{\textwidth}{>{\raggedright\arraybackslash}p{3.4cm} Y}
\rowcolor{AccentDark}
\thead{Component} & \thead{What it means} \\
\toprule
Self-awareness & Awareness of your own emotions (e.g., recognizing ``I am getting angry''). Built by habitual introspection and meditation. \\
Social awareness & Recognizing and understanding the emotions of others. \\
Self-regulation & Managing and controlling your own emotions (e.g., anger management). \\
Social regulation & Relationship management --- influencing/regulating others' emotions; managing relationships. \\
Motivation & The self-motivation that ultimately drives regulation of emotions and delayed gratification. \\
\bottomrule
\end{tabularx}
\end{center}

\paragraph{Why self-awareness matters} Until we have awareness about ourselves, how will we control our emotions? Meditation simply asks you to focus on yourself --- this is what the Buddha did. So self-awareness is very important; it is the first step to processing and regulating emotions.

\paragraph{Why social awareness matters} The second step is social awareness --- understanding the emotions of others --- because if I do not understand the other's situation, how will I do relationship management? Several examples illustrate this:

\paragraph{Tribals --- two approaches and the golden mean} There were two choices for dealing with tribals. One is the ``museum-species approach'' --- leaving tribals as they are, like specimens kept untouched in a museum, saying in effect ``keep living in the forest.'' The other extreme is full assimilation, where the tribal's own culture, language and traditions are lost as they adopt the mainstream language and culture. The Nehruvian Panchsheel --- the Tribal Panchsheel policy --- is the golden mean between these two approaches: not isolation, and not forced assimilation, but ``I should help them while keeping their culture, language, etc.\ intact.'' This became possible because of better awareness.

\begin{GovtPolicy}
\begin{itemize}
\item Nehruvian / Tribal Panchsheel: the golden mean between the ``museum-species'' (isolation) approach and forced assimilation.
\item Principle: help tribals develop while keeping their culture and language intact --- enabled by better social awareness.
\end{itemize}
\end{GovtPolicy}

\paragraph{Naxalism --- security vs.\ development, and SAMADHAN} Naxalism can be seen purely as a law-and-order problem (send in security forces, a security-oriented approach) or through a ``development-deficit'' and humanitarian lens. Better social awareness is needed to diagnose the problem correctly. The SAMADHAN approach is a combination of both --- it became possible with better social awareness. Without better social awareness, a problem cannot be diagnosed correctly.

\begin{GovtPolicy}
\begin{itemize}
\item SAMADHAN: the Government's strategy against Naxalism --- a combination of the security approach and the development-deficit/humanitarian approach.
\item Correct diagnosis of Naxalism requires better social awareness.
\end{itemize}
\end{GovtPolicy}

\paragraph{Swachh Bharat --- building toilets is not enough} Merely building toilets does not automatically bring cleanliness or reduce open defecation. Why not? Because of poor social awareness. Correct social awareness recognizes that Swachh Bharat Mission requires behavioural change, not just infrastructure. In all these problems, without better social awareness we cannot solve the problem --- which shows the importance of emotional intelligence.

\begin{ExampleBox}
\begin{itemize}
\item Toilets built $\neq$ cleanliness achieved, unless behaviour changes --- requires social awareness.
\item Beti Bachao Beti Padhao $\rightarrow$ change attitude towards the girl child.
\item Swachh Bharat $\rightarrow$ not just build toilets but change behaviour (``use them too'').
\end{itemize}
\end{ExampleBox}

\paragraph{Non-verbal communication and power distance} Social awareness is also important because many times a person cannot state their problem in words. In a live class, the teacher reads students' faces to judge whether they have understood --- this is non-verbal communication, which is impossible without social awareness. An officer must read facial expressions and body language to understand how much difficulty the other person is facing. This matters because there is a lot of power distance between officers and the people; officers must themselves have social awareness to better regulate their behaviour.

\begin{KeyConcept}
\begin{itemize}
\item Non-verbal communication (reading faces, expressions) depends on social awareness.
\item Large power distance between officers and people makes officers' social awareness essential.
\end{itemize}
\end{KeyConcept}

\paragraph{Self-regulation} Self-regulation means managing and regulating your own emotions. The Buddha's example: he was surrounded by disciples, and a person on the way spat on him; through his self-regulated, calm reaction the person was humbled (``paani-paani ho gaya'') and converted to become his disciple. That is self-regulation --- managing and regulating your emotion.

Nelson Mandela's forgiveness is another example. When Mandela came out of jail, he too had to practise self-regulation, because anger is natural when injustice has been done to you. ``Forgiveness is a tribute of the strong; the weak cannot forgive.'' To forgive, your heart has to be a little bigger --- you must be able to digest your anger.

Also: ``Anger and prejudice are enemies of correct understanding'' (a UPSC-relevant line). If we are angry, our understanding is not as good; the mind, when agitated or angry, does not reflect things clearly. So anger and intolerance are enemies of correct understanding.

\begin{PYQLink}
\begin{itemize}
\item ``Anger and intolerance are enemies of correct understanding.'' (quotation theme)
\item Self-regulation examples: the Buddha (spat upon, calm response); Nelson Mandela's forgiveness; ``forgiveness is a tribute of the strong''.
\end{itemize}
\end{PYQLink}

\paragraph{Social regulation (Relationship management)} Social regulation is relationship management --- managing others' emotions, giving people full attention. When someone is experiencing grief, how do I manage my relationship with that person? This is where the IQ-versus-EQ discussion becomes relevant.

\paragraph{IQ vs.\ EQ} Traditionally we relied heavily on IQ (intelligence quotient) --- the belief that if you have high IQ you will be successful. But in real life, emotional quotient (EQ) has a very big role. There are many people who are intellectually sound with good IQ, but when placed in a team or asked to lead one, they cannot manage people, because they lack EQ. EQ is your ability to understand others' emotions and regulate emotions.

\begin{center}
\begin{tabularx}{\textwidth}{Y Y}
\rowcolor{AccentDark}
\thead{IQ (Intelligence Quotient)} & \thead{EQ (Emotional Quotient)} \\
\toprule
Reasoning, problem-solving, cognitive ability. & Emotional skills --- understanding and managing emotions of self and others. \\
Traditionally over-relied upon. & Crucial for leadership, team management and relationship management. \\
Measured by conventional IQ tests. & About making emotions ``work for you rather than against you''. \\
\bottomrule
\end{tabularx}
\end{center}

For leadership, EQ is essential. A person with high IQ but poor EQ cannot deliver their best as a team leader. Respect is commanded by following EQ: leaders like Gandhiji became great leaders because they had EQ --- people relate to you when you understand their pain and emotions, and only then do you command respect. A leader's big task is to hold a vision and a goal; the leader's goal becomes the followers' goal only when the leader convinces them. Gandhiji's goal was Swaraj/independence, and he made everyone understand why we all should follow it. Without EQ you cannot persuade others; ``1 + 1 = 11'' --- the synergy --- comes only with EQ. For a civil servant (IAS/IFS), EQ is very important for relationship management.

\begin{ExampleBox}
\begin{itemize}
\item Leadership examples of EQ: Gandhiji (persuaded masses towards Swaraj); a leader's goal becomes followers' goal only through EQ.
\item High IQ + poor EQ $\rightarrow$ cannot lead a team effectively.
\end{itemize}
\end{ExampleBox}

\paragraph{Motivation \& the Marshmallow Test} To manage emotions, you ultimately need to be motivated. The Marshmallow Test (a US university experiment): children were seated at a table with a simple rule --- they could leave anytime and get one marshmallow, but if they waited (about an hour) they would get two marshmallows. Some children ate the one marshmallow immediately; others delayed. On studying the group, those who were able to delay their gratification turned out to be much more successful in life, because they were motivated enough to delay gratification.

Delayed gratification means you are motivated enough to postpone immediate satisfaction for a bigger reward --- the fun of a party today can be set aside because ``I want to become an IAS or a public servant, so I delay this gratification and take some pain now.'' Many people leave a job or a bigger campus-placement package to prepare for something bigger --- that too is delaying gratification. Motivation is what enables you to regulate your emotions and delay gratification. Even a bureaucrat must delay gratification and motivate --- working for what is right rather than merely for postings and promotions.

\begin{DataFact}
Marshmallow Test (US university): 1 marshmallow now vs.\ 2 if you wait $\sim$1 hour. Children who delayed gratification were much more successful later --- due to motivation.
\end{DataFact}

\begin{KeyConcept}
\begin{itemize}
\item Motivation $\rightarrow$ enables delayed gratification $\rightarrow$ enables emotional regulation.
\item Applies to civil servants: work for what is right, not merely for postings/promotions.
\end{itemize}
\end{KeyConcept}

\subsection{Case Study: The Bank Manager \& the Digital Arrest}

A person (the victim) repeatedly came to a bank wanting to withdraw large sums. The whole case was one of ``digital arrest'' --- a scam. This is a great example of emotional intelligence:

\begin{itemize}
\item Social awareness: No one told the bank manager that the customer was under digital arrest; the manager herself understood, from the customer's behaviour, that she was in trouble.
\item Self-regulation: When the manager intervened, the customer replied that she should ``mind her own business.'' Anyone would be angered by such a reply, but the manager regulated her emotions.
\item Compassion \& relationship management: The manager still pursued the matter and spoke with her, ultimately saving her from losing about \rupee 79 lakh.
\end{itemize}

As written up: the victim returned to dissolve another savings worth \rupee 30 lakh; this time the manager refused to transmit any money unless someone from the family arrived. By this time the victim was agitated and scared, and on the steps of the bank was pleading with someone on the phone. The manager, watching her movements on CCTV, walked out and took the phone. The person on the other end claimed to be the victim's brother-in-law; when asked the sister's name he fell silent, and on continued probing, disconnected the call.

\begin{ExampleBox}
\begin{itemize}
\item Bank manager (digital-arrest case): social awareness (sensed trouble unaided) + self-regulation (``mind your own business'' reply) + compassion + relationship management.
\item Outcome: saved the victim $\sim$\rupee 79 lakh; caught the scam via CCTV and probing the caller.
\item Shows the power of non-verbal communication and proactiveness.
\end{itemize}
\end{ExampleBox}

\subsection{Case Study: Topography of Terror (Berlin)}

The ``Topography of Terror'' and the ``Memorial to the Murdered Jews of Europe'' are in Berlin. In the museum there are many letters displayed, which visitors are asked to read (not merely watch). At a site where the police headquarters once stood, officials discussed how to kill Jews on an ``industrial scale.'' The reasoning recorded there was that if Jews were killed one by one with guns, the killers would suffer negative mental-health consequences --- so a method was devised to kill people on an industrial scale, since they intended to kill lakhs of people.

One method recorded: prisoners were placed in the windowless storage area of a vehicle, and during the journey to a nearby forest, engine exhaust fumes were pumped into that sealed area; the prisoners inside died an agonizing death, and their bodies were buried in mass graves in the forest. Visitors, on reading these accounts, feel the pain of others. On leaving the ``Topography of Terror,'' visitors encounter reflective questions asking them to introspect, making them aware of historical blunders and prompting an apologetic reckoning.

\begin{ExampleBox}
\begin{itemize}
\item Topography of Terror / Memorial to the Murdered Jews of Europe (Berlin): museums sensitize visitors, building social awareness.
\item Nazis planned killing on an ``industrial scale'' partly to spare killers' mental health --- a chilling case of misused rationality.
\item Museums use historical example to sensitize people (build EI) about victims' plight.
\end{itemize}
\end{ExampleBox}

\begin{KeyConcept}
You need emotional intelligence to experience compassion --- EI is the ability that lets you feel others' pain; compassion and empathy follow from it.
\end{KeyConcept}

\subsection{More Examples of Emotional Intelligence (and its absence)}

\paragraph{Lack of emotional intelligence}
\begin{itemize}
\item BLO overburdening: State governments were not empathetic to their own workforce; instead of deploying fresh staff to ease each BLO's workload, the burden was left on them, and many BLOs committed suicide because they were overburdened. The Supreme Court remarked on this. A lack of emotional intelligence.
\item Officials slapping citizens: A person coming out of a Covid queue was slapped; in a Rajasthan petrol-pump incident an SDM reportedly slapped an attendant. These show bureaucrats unable to manage their emotions.
\end{itemize}

\paragraph{Presence of emotional intelligence}
\begin{itemize}
\item Chadotaru tradition (Banaskantha): The Banaskantha Police Station initiated a meticulous process of reconciliation to ensure a peaceful homecoming for families who had had to leave their village. The police understood the community's tradition (Chadotaru), regulated emotions, and corrected a historical error --- bringing 29 (and separately 21) tribal families back. A great example of emotional intelligence (social awareness + relationship management).
\item Ganguly's leadership: Players like Yuvraj Singh and Harbhajan Singh said they ``could die for a captain like Ganguly.'' Virender Sehwag, Harbhajan and Yuvraj say Ganguly backed them fully. Leadership means backing your people and understanding their emotions, so a new player does not feel anxiety about being dropped after one failure --- anxiety harms performance.
\item Rahul Dravid's press conferences: When India lost the 2023 World Cup final, Rahul Dravid came and took the press conference; when India won the 2024 T20 World Cup, Rohit Sharma took the press conference. Dravid took the blame in defeat --- a great example of leadership and EI, because he knew what the players were going through.
\item Sangi Express (bike ambulance): IAS officer Awanish Sharan (Chhattisgarh) started a bike ambulance named ``Sangi Express'' because in remote areas regular ambulances could not reach and were costly. This low-cost two-wheeler ambulance reduced medical costs. Example of EI + innovation + compassion.
\item SDM carrying the injured: When a couple lay bleeding on the road, an SDM took them to hospital --- by car and even on a passer-by's bike --- understanding and trying to elevate their pain. An example of EI and compassion.
\item Osmanabad Collector: The Osmanabad Collector sat on the ground floor to meet a person with a disability (a person with different ability) who could not sit on a chair. The visitor did not ask the officer to sit on the ground; the Collector did it on his own to understand and hear them. An example of EI.
\item PM waiting at the airport: The PM waiting at the airport for some time before an exam to avoid causing traffic --- an example of EI and compassion (anticipating others' pain).
\item Liquor ban in religious cities (MP): Liquor ban in 19 religious cities in Madhya Pradesh (under select gram panchayats) so that people's religious sentiments are not hurt --- an example of social awareness / respecting cultural traditions.
\end{itemize}

\begin{ExampleBox}
\begin{itemize}
\item EI present: Chadotaru reconciliation (Banaskantha Police); Ganguly \& Dravid leadership; Sangi Express bike ambulance (Awanish Sharan); Osmanabad Collector; PM waiting at airport; MP liquor ban in 19 religious cities.
\item EI absent: BLOs overburdened (suicides; SC remarked); officials slapping citizens (Covid queue, Rajasthan petrol pump).
\end{itemize}
\end{ExampleBox}

\begin{KeyConcept}
All examples of compassion/empathy can also be used as examples of emotional intelligence --- they are interchangeable (e.g., Kiran Bedi's Tihar Jail reforms understanding prisoners' plight; welfare schemes).
\end{KeyConcept}

\begin{QuickRevision}
\begin{itemize}
\item EI = ability to interpret, understand \& manage one's own and others' emotions (Recognition + Regulation).
\item 5 components: self-awareness, social awareness, self-regulation, social regulation (relationship management), motivation.
\item Being an ability, EI can be MISUSED and does not by itself make one moral.
\item Golden-mean policy examples: Tribal Panchsheel; SAMADHAN (Naxalism).
\end{itemize}
\end{QuickRevision}

\section{Trained Apathy \& Functional Detachment}

This topic addresses whether too much sensitivity can itself be a problem. A public servant conducting a press conference on a very sensitive matter while grinning (``Poem case: cop caught grinning during media briefing, video sparks fury online'') is not in alignment with what is expected --- showing why emotional intelligence is important.

Consider a doctor who deals with patients all day, each going through some problem. Of course the doctor should be sensitive; but if that sensitivity goes to the extent that the doctor breaks down and cries for every cancer patient, imagine the emotional toll. Similarly, imagine people working in a mortuary. There is a functional perspective to apathy: a certain apathy/detachment helps them do their work properly. The point is not that one should become apathetic, but that too much sensitivity towards emotions can also be harmful, because you are doing stressful work all day.

When you become a DM or SP, many people will come to you with requests and problems; you should be sensitive, but if you become emotional every time and, say, cry along with everyone, your own mental health will suffer. This is why a certain element of detachment is necessary. Emotional intelligence means controlling your emotions --- being emotional is not the same as being emotionally intelligent. As the quote says, anyone can be angry, but to be angry with the right degree, measure, person and time is difficult. We should feel emotions but keep our own emotions in check, and not become excessively emotional, because that creates problems.

\begin{KeyConcept}
\begin{itemize}
\item ``Trained apathy'' / functional detachment: a measured detachment helps professionals (doctors, mortuary staff, DMs/SPs) function amid constant stress.
\item Being emotional $\neq$ being emotionally intelligent; EI = keeping one's own emotions in check.
\end{itemize}
\end{KeyConcept}

\begin{CriticalPoint}
Too much sensitivity can harm the officer's own mental health --- a certain element of detachment is necessary.
\end{CriticalPoint}

\section{Using Emotions Positively (Emotions Working For You)}

Emotional intelligence is about making emotions work for us rather than against us. ``How's the Josh?'' is an example of elevating emotions and channelling them in the right way. The key case study is the use of non-rational motives to change attitudes.

\subsection{Rational vs.\ Non-rational (Extra-rational) Behaviour}

If we assume every human being is a rational human being (one who thinks logically and objectively), then giving them new information should make them act on it. But behaviour is not always guided by rational variables. For example, a smoker has been told ``smoking causes cancer,'' yet still smokes --- that behaviour is not rational. This means actual behaviour is determined by extra-rational / non-rational factors too, i.e., emotions (for instance, addiction). Similarly, open defecation is not addressed merely by logically explaining the diseases it causes, because human beings do not always operate logically; they operate emotionally too. This is why emotions are used to change people's attitudes.

\begin{KeyConcept}
\begin{itemize}
\item Actual behaviour is driven by extra-rational (non-rational) factors --- emotions --- not logic alone.
\item Logical motive (e.g., ``I won't smoke because it causes cancer'') vs.\ non-rational motive (smoking despite knowing).
\end{itemize}
\end{KeyConcept}

\subsection{Swachh Bharat Mission --- Pride, Disgust, Love}

Swachh Bharat Mission used three important non-rational motives --- pride, disgust and love --- to change attitudes.

\paragraph{Disgust} In workshops, a facilitator would drop a dead fly into a glass of water, take it out, and offer the water to drink. Naturally no one drinks it, because an element of disgust is present. The facilitators used this to explain that when a fly sits on human faeces and then lands on food, people still eat that food. Advertisements show this pictorially to evoke disgust, because people are not reasoning logically about disease --- the disgust emotion is used to change attitudes.

\paragraph{Pride} Big celebrities were involved in Swachh Bharat Mission, wielding brooms, to convey that cleaning work is something to take pride in. If even prominent people are cleaning their surroundings, ordinary citizens can take pride in doing the same. Celebrities posted videos of cleaning their neighbourhoods to evoke pride.

\paragraph{Love} Marriages were reportedly cancelled where the groom's house had no toilet --- invoking the element of love: if you truly love your future wife/daughter-in-law, you would build a toilet. Where there is a lack of toilets, women going outside for defecation face greater risk of crime and safety problems; so the love element (``care for your daughters' and daughters-in-law's safety'') was invoked to encourage toilet construction.

\begin{GovtPolicy}
\begin{itemize}
\item Swachh Bharat Mission used three non-rational motives --- PRIDE, DISGUST, LOVE --- to drive attitudinal change.
\item Disgust: dead-fly-in-water demonstration; pride: celebrities with brooms; love: no-toilet marriage cancellations \& women's safety.
\end{itemize}
\end{GovtPolicy}

\subsection{Pictorial Warnings on Tobacco}

On tobacco products, the text ``smoking causes cancer'' is information addressed to logic. But the graphic photo is added to evoke your emotions. A significant portion --- reportedly around 40\% --- of the pack should carry the warning, precisely to evoke emotions. Both this and the Swachh Bharat case are strong case studies of ``using emotions for you rather than against you.'' (The film character Mukesh, shown in Indian cinema anti-tobacco clips, is a familiar reference.)

\begin{DataFact}
\begin{itemize}
\item Pictorial warnings must cover a significant portion (reportedly $\sim$40\%) of tobacco packaging to evoke emotion.
\item Text warning = information (logic); graphic image = emotional appeal.
\end{itemize}
\end{DataFact}

\begin{PYQLink}
\begin{itemize}
\item PYQ theme: ``Emotional intelligence is the ability to make your emotions work for you instead of against you. Do you agree with this view?''
\item PYQ theme: ``Apart from intellectual competency and moral qualities, empathy and compassion are vital attributes\ldots\ that help the civil servant tackle crucial issues.''
\end{itemize}
\end{PYQLink}

\begin{QuickRevision}
\begin{itemize}
\item Positive use of emotions: SBM (pride, disgust, love); tobacco pictorial warnings ($\sim$40\%).
\item Negative examples: officials slapping citizens; the ``grinning cop'' press conference.
\item Behaviour is driven by extra-rational factors $\rightarrow$ attitudes are changed via emotions.
\end{itemize}
\end{QuickRevision}

\section{Can Emotional Intelligence Be Learned?}

Yes --- emotional intelligence can be learned and inculcated. There are many methods and programmes, because the key components of EI (understanding and regulating one's emotions) can be developed.

\subsection{Methods to Inculcate Emotional Intelligence}

\begin{itemize}
\item Training programmes / workshops / sensitivity training: Many courses include sensitivity training to sensitize people about others' situations and problems.
\item Role playing: For example, an actor preparing to play a transgender role spends many days meeting transgender persons to enter that role and understand their perspective. Role-playing exercises help build EI. (A husband who does household work himself --- previously dismissed as trivial --- realises its difficulty; a good role-playing example.)
\item Role modelling / mentorship / learning from leaders: e.g., M.S. Dhoni mentoring someone on captaincy would stress processing one's emotions and staying calm in stressful situations; learning from the Buddha's calm reaction. Learning from leaders is very important.
\item Enhancing self-awareness: by habitual (self-)introspection and meditation.
\item Enhancing social awareness: e.g., LBSNAA's ``Bharat Darshan,'' where officer-trainees are taken to different areas to understand the problems of different regions; Gandhiji's train journey across India to build social awareness; media, which can sensitize people (``what the eyes do not see, the heart does not grieve about'' --- e.g., the Manipur incident used to spread awareness).
\item Regulating emotions: yoga, meditation, pranayama, deep breathing (e.g., the deep breaths candidates take before a UPSC interview), and diversion of activity (as Virender Sehwag reportedly diverted attention --- singing a song, punching a pillow --- to divert oneself).
\item Religious lessons: e.g., Buddhism's focus on Karuna (compassion) --- understanding others' pain and removing it; every religion teaches such lessons.
\end{itemize}

\begin{GovtPolicy}
\begin{itemize}
\item LBSNAA ``Bharat Darshan'' builds social awareness among officer-trainees by exposing them to different regions.
\item Media can inculcate EI by sensitizing people (e.g., Manipur incident); ``what the eyes do not see, the heart does not grieve about''.
\end{itemize}
\end{GovtPolicy}

\begin{KeyConcept}
\begin{itemize}
\item Self-awareness $\rightarrow$ habitual introspection + meditation.
\item Social awareness $\rightarrow$ tours (Bharat Darshan), media, role modelling.
\item Regulation $\rightarrow$ yoga, meditation, pranayama, deep breathing, diversion of activity.
\item Also: sensitivity training, role playing, mentorship/learning from leaders, religious lessons.
\end{itemize}
\end{KeyConcept}

\subsection{Emotional Intelligence Can Be Misused}

Because EI is only an ability, it can be misused. Hitler was emotionally intelligent --- what we call relationship management can, for someone with unethical character, become manipulation. EI means understanding others' pain and managing (motivating) relationships; but if your character is unethical, you use that understanding to change others' attitudes for your own benefit --- which is manipulation. Hitler convinced all of Germany that, ``for the rebirth of Germany, killing of Jews is important,'' and most stayed silent while Jews were killed.

Political leaders and religious leaders sometimes mislead their followers for their own benefit; they are emotionally intelligent --- they understand what statement will produce what reaction --- but their intentions are unethical, so they misuse their ability to manage emotions. Therefore emotional intelligence, by itself, does not lead to ethical decision-making, because the ability can be misused.

\begin{CriticalPoint}
\begin{itemize}
\item EI is value-neutral (an ability) $\rightarrow$ it can be misused; by itself it does NOT ensure ethical decision-making.
\item Misuse examples: Hitler; hate speeches by political leaders; manipulation by some religious leaders; terrorist brainwashing / ISIS recruitment \& propaganda; misleading advertisements (e.g., ``fair and lovely'' type ads); ``divide and rule'' (British).
\end{itemize}
\end{CriticalPoint}

\begin{ExampleBox}
Media is a double-edged tool: it can sensitize (build EI) OR be misused (misleading ads, propaganda).
\end{ExampleBox}

\begin{QuickRevision}
\begin{itemize}
\item EI CAN be learned: training, role playing, role modelling/mentorship, introspection, meditation, yoga/pranayama/deep breathing, diversion, religious lessons, media.
\item Enhance self-awareness by habitual introspection; social awareness by Bharat Darshan / media.
\item EI can be MISUSED (Hitler, hate speech, ISIS propaganda) $\rightarrow$ EI alone $\neq$ ethical action.
\end{itemize}
\end{QuickRevision}

\section{Attitude: Meaning \& Structure}

\subsection{What is Attitude?}

Attitude is a relational term. ``Relational'' means that attitude is always discussed with respect to something; you never ask about attitude ``in general,'' but always attitude towards a particular attitudinal object.

Definition: Attitude is your predisposition (or tendency) to act in a certain way towards an attitudinal object.

\begin{KeyConcept}
\begin{itemize}
\item Attitude = predisposition / tendency to act in a certain way towards an attitudinal object.
\item Attitude is always RELATIONAL --- measured only with respect to an attitudinal object.
\end{itemize}
\end{KeyConcept}

An attitudinal object can be anything: a person, a song, a singer, an actor, a particular ideology, a food item, or any race, caste, sex or religion, etc. For example, my predisposition towards a particular singer or song shows whether my attitude is positive or negative towards it; my attitude towards a particular political party (say, positive or negative) is a political attitude. Predisposition is the tendency --- e.g., I like Arijit Singh's songs, so my tendency is to listen to them (positive attitude); I do not enjoy a certain actor's acting, so my predisposition is not to watch that actor's movies. Likes and dislikes about food (mushroom, paneer, eggs) are likewise attitudes towards those objects.

\begin{ExampleBox}
\begin{itemize}
\item Attitudinal objects can be: a person, singer, actor, song, ideology, political party, food, or a race/caste/sex/religion.
\item Positive attitude $\rightarrow$ tendency to approach (e.g., listen to a liked singer); negative attitude $\rightarrow$ tendency to avoid.
\end{itemize}
\end{ExampleBox}

\subsection{Structure of Attitude --- the CAB (ABC) Model}

Every attitude has three major components. The model is called the CAB (or ABC) model:

\begin{center}
\begin{tabularx}{\textwidth}{>{\raggedright\arraybackslash}p{2.4cm} >{\raggedright\arraybackslash}p{4.3cm} Y}
\rowcolor{AccentDark}
\thead{Component} & \thead{Meaning} & \thead{Illustration} \\
\toprule
Cognitive & Your belief pattern --- whatever you hold to be true; the information you have. This shapes your attitude. & ``Trees conserve the environment, stop soil erosion, aid wildlife, give food \& fibre'' $\rightarrow$ belief supporting a positive attitude to trees. \\
Affective & How you FEEL about the attitudinal object --- your emotional overtone. & Liking greenery / feeling happy on seeing trees; feeling disgust at a cockroach. \\
Behavioural & Your behavioural tendency --- what you would actually do. & A pessimist tends to run from challenges; a thief tends to run on seeing police; if you can, you plant a tree. \\
\bottomrule
\end{tabularx}
\end{center}

\paragraph{Worked example --- attitude towards trees / tree plantation} Cognitive (belief): trees conserve the environment, stop soil erosion, support wildlife conservation, and provide food and fibre --- so I believe trees are beneficial. Affective (feeling): I like greenery and feel happy seeing trees around me. Behavioural (tendency): if I get the chance, I will plant trees and motivate others to do so. All three components correspond to the same object (trees).

\paragraph{Worked example --- racial prejudice (for illustration only)} Cognitive: a belief that ``whites are intellectually superior.'' Affective: feeling hatred/discomfort on seeing a black person. Behavioural: a tendency to discriminate. This combination is what constitutes prejudice --- a belief, an emotional overtone, and a discriminatory behavioural tendency towards the object. (Used purely to illustrate the components, not as the teacher's view.)

\begin{KeyConcept}
\begin{itemize}
\item CAB model: Cognitive (belief) + Affective (feeling/emotion) + Behavioural (behavioural tendency).
\item Prejudice = belief + hatred (affect) + discriminatory tendency towards an object.
\end{itemize}
\end{KeyConcept}

\subsection{Components May Not Be in Congruence}

The three components are not always in sync. Sometimes one dominates the others:

\begin{itemize}
\item Bitter medicine: Cognitive is strong (I believe it will treat me), affective is very negative (it is bitter), yet the behavioural outcome is that I take it --- the cognitive component dominates. Generally, when something is very important in our life, the cognitive component usually dominates.
\item Cheap impulse purchase: You know it is wasteful (cognitive: ``I don't need it''), but you like it (affective is strong) and it is cheap, so you buy it --- the affective component dominates. In consumerism, smoking addiction and drug addiction, the affective component tends to dominate even when we know the harm.
\item Congruence: Sometimes all three align --- e.g., I believe I should exercise (cognitive), I enjoy the gym (affective), and I actually go (behavioural).
\end{itemize}

\begin{CriticalPoint}
\begin{itemize}
\item The three components (C, A, B) may or may not be congruent.
\item For important matters $\rightarrow$ cognitive usually dominates (e.g., bitter medicine).
\item For consumerism/addictions $\rightarrow$ affective usually dominates (impulse buys, smoking, drugs).
\end{itemize}
\end{CriticalPoint}

\begin{PYQLink}
Gandhiji's statement (quotation theme): ``What you think, you become'' --- belief pattern shapes attitude, which shapes behaviour and ultimately character.
\end{PYQLink}

\section{Functions of Attitude}

``Function'' means the utility of something for you --- how does an attitude help you in your life? (This functional perspective is familiar from anthropology/sociology.) There are four important functions:

\subsection{Knowledge Function}

Attitude helps us make sense of the outside world without having to evaluate every object again and again. Think of an attitude as a ``scheme'' or classification --- a pair of glasses through which we categorize things.

\begin{itemize}
\item Snakes: We know snakes can kill, so our usual attitude towards snakes is cautious/fearful --- we keep a distance. Even though 60--70\% of snakes are non-venomous, we still keep our distance, because it saves our life; we cannot verify each snake, and it might turn out to be venomous.
\item Water: If water looks muddy or smells off, we do not drink it --- even if it might be clinically safe --- because our attitude is that clean water is odourless and clear. We do not clinically test every glass of water we ever drink; the attitude simplifies life.
\item Medicine: We take medicine without evaluating whether it is fake each time, because we have formed the attitude that government-checked medicine is fine.
\end{itemize}

So the knowledge function reduces cognitive load and simplifies life. Its negative consequence is stereotyping --- e.g., assuming all snakes are dangerous, or that all trees are good (whereas some, like the eucalyptus tree, are water-intensive and problematic where water is scarce).

\begin{KeyConcept}
\begin{itemize}
\item Knowledge function: attitude acts as a ``lens/scheme'' to make sense of the world and reduce cognitive load --- no need to re-evaluate every object.
\item Negative side: STEREOTYPING (e.g., ``all snakes dangerous''; ``all trees good'' ignoring water-intensive eucalyptus).
\end{itemize}
\end{KeyConcept}

\subsection{Utilitarian / Socially-Adjustive Function}

Some attitudes help us adjust to, and be accepted by, our social group, by conforming to expected norms. We adopt certain attitudes to be likeable and socially acceptable within a group (relevant, for instance, to how politicians present themselves to groups and at elections). This is the socially-adjustive (utilitarian) function --- the attitude helps you adjust to or be welcomed by the group because you are conforming to expected norms. (This also links to why honesty is easier to sustain --- a truthful attitude spares you the tangle of remembering and defending lies later.)

\begin{KeyConcept}
Socially-adjustive / utilitarian function: attitudes that help you fit in and be accepted by conforming to group norms.
\end{KeyConcept}

\subsection{Ego-Defensive Function}

Some attitudes are held to protect our own mental health and self-image. Adopting a certain attitude helps us defend ourselves against criticism. For instance, the doctor's measured detachment helps him defend his own mental health; or accepting in advance that ``some people will always criticize me, whatever I do'' protects one from being wounded by that criticism (even Lord Ram was questioned by people). This is the ego-defensive function.

\begin{KeyConcept}
Ego-defensive function: attitudes adopted to protect one's own mental health and self-image against criticism.
\end{KeyConcept}

\subsection{Value-Expressive Function}

Some attitudes we ``wear on our sleeve,'' because they express who we are --- our identity and character. Virat Kohli's aggressiveness and Dhoni's calmness are attitudes they carried because these express their identity. Another example: in inter-caste relationships, someone may say ``my father will never agree,'' because the attitude against inter-caste marriage is very important to that parent's identity --- they disapprove of it entirely, so the attitude expresses their identity and is very significant to them. Danveer Karna (his commitment to donation) and Bhishma Pitamah (his vow) are cases where value-expressiveness was extremely important.

\begin{KeyConcept}
Value-expressive function: attitudes that express one's identity/character (e.g., Kohli's aggression, Dhoni's calm; Danveer Karna; Bhishma).
\end{KeyConcept}

\begin{QuickRevision}
Four functions of attitude: (1) Knowledge (lens; reduces cognitive load; risks stereotyping); (2) Utilitarian/Socially-adjustive (fitting in); (3) Ego-defensive (protects mental health/self-image); (4) Value-expressive (expresses identity).
\end{QuickRevision}

\section{Attitude's Influence on Thought \& Behaviour}

\subsection{Influence on Thought}

There is a cyclical relationship. The cognitive component (belief) goes in as an input into attitude formation. But once an attitude has been formed, it in turn influences all the subsequent thoughts that arise in your mind.

Example: if I have already formed a casteist attitude (from a belief that certain ``menial/polluting'' jobs should be done by certain castes), then when a social incident occurs and WhatsApp forwards reach me, I will most likely fall prey to confirmation bias. Confirmation bias means I treat as true the information that confirms my existing attitude/bias --- ``I already knew people of that caste behave this way,'' even if the forward is false, because my attitude was already formed.

Another example: if I have a communal attitude and receive information that ``that community did something wrong,'' I will most likely accept it as true even if it is fake, because my attitude is already set. So attitude shapes our thought process: thoughts (beliefs) form the attitude, and the attitude then shapes further thoughts.

\begin{KeyConcept}
\begin{itemize}
\item Cyclical: belief (cognitive) $\rightarrow$ attitude formed $\rightarrow$ attitude then filters/shapes further thoughts.
\item Attitude $\rightarrow$ confirmation bias: we accept information that confirms our existing attitude, even against evidence.
\end{itemize}
\end{KeyConcept}

The patriarchy example: if a person holds a patriarchal attitude and a rape incident occurs, he thinks ``I already knew --- they roam in short clothes, so this happens,'' even when there is contradictory evidence (e.g., the victim was fully/traditionally clothed). NCRB data show that most sexual-harassment cases involve people known to the victim --- but the biased person will ignore this and see only what fits the attitude. ``As is the lens on the eye, so is the vision.''

\begin{DataFact}
Per NCRB data, most sexual-harassment cases involve persons known to the victim --- evidence a biased attitude tends to ignore.
\end{DataFact}

\begin{PYQLink}
\begin{itemize}
\item Quotation theme: ``We don't see things as they are; we see things as we are.'' (We view the object through our own lens/attitude.)
\item Related: ``Your perception of me is a reflection of you; my reaction to you is an awareness of me.''
\end{itemize}
\end{PYQLink}

\subsection{Influence on Behaviour}

Attitude is only a predisposition/tendency; it may or may not translate into actual behaviour. This line is very important: an attitude may or may not translate into actual behaviour.

\begin{itemize}
\item Mismatch example: Many politicians may hold a patriarchal attitude but not display it when meeting people or giving speeches; their public behaviour may keep a full distance from their actual attitude. A bureaucrat may speak against corruption publicly while holding an attitude that is quite positive towards it.
\item Everyday example: You may have a positive attitude towards study (you feel like studying) but not actually study --- so a positive attitude alone does not let us predict the behaviour.
\item Hypocrisy: e.g., a person who plants a tree but also cuts trees; a police officer who beats his wife in private life. Attitude and behaviour need not be congruent.
\end{itemize}

\begin{CriticalPoint}
\begin{itemize}
\item Attitude = tendency/predisposition; it may or may not translate into actual behaviour.
\item Hence attitude--behaviour mismatch (e.g., anti-corruption speeches vs.\ corrupt attitude; positive study attitude but no studying) --- and hypocrisy.
\end{itemize}
\end{CriticalPoint}

\paragraph{Factors determining when attitude predicts behaviour}

\begin{center}
\begin{tabularx}{\textwidth}{>{\raggedright\arraybackslash}p{3.4cm} Y}
\rowcolor{AccentDark}
\thead{Factor} & \thead{How it strengthens the attitude--behaviour link} \\
\toprule
Strength of attitude & A stronger attitude is more likely to be in alignment with behaviour. E.g., Gandhiji's very strong attitude to non-violence lets us predict he would not support violence; Danveer Karna's strong attitude to donation predicts he would never refuse to donate. \\
Specificity of attitude & The more specific the attitude, the more strongly it guides behaviour. ``I believe in sustainability'' is broad; ``I have a negative attitude towards plastic straws'' is specific and guides behaviour better. ``Study 8 hours daily'' is more specific (and predictive) than ``study a lot''. \\
Accessibility of attitude & Attitudes we think about frequently are more accessible and more likely to guide behaviour (e.g., someone constantly thinking about tree plantation is more likely to plant/motivate others). \\
Moderating role of external factors & The surrounding environment can change behaviour --- e.g., an honest person in a very corrupt office; whether the group is liberal or not. Politicians don't reveal everything publicly because the surrounding factors are not in their favour. \\
Perceived behavioural control (Theory of Planned Behaviour) & If you feel you can actually carry out the behaviour in the circumstances (control over your own behaviour), you are more likely to act; if not, you stay silent or comply. \\
\bottomrule
\end{tabularx}
\end{center}

\paragraph{Perceived behavioural control --- examples} An honest person joining a very corrupt office may hold a strong, specific, accessible positive attitude (``I will not take a bribe''), but external factors and perceived control matter: if being honest could get him harmed, his perceived control is low, so he may stay silent, or show indifference to bribery (``take it if you must; I won't''), without openly protesting. Similarly, a girl may hold a strongly negative attitude towards dowry, yet dowry is still given at her marriage, because she perceives little control --- protesting might mean even her own parents will not listen. Perceived behavioural control is the extension (of the model) that captures how effectively you feel you can perform the behaviour in the given circumstances.

\begin{KeyConcept}
\begin{itemize}
\item Theory of Planned Behaviour: perceived behavioural control --- how effectively you believe you can perform the behaviour in the circumstances --- strongly affects whether attitude becomes action.
\item Low perceived control (e.g., honest officer in corrupt office; girl vs.\ dowry) $\rightarrow$ attitude may not translate into behaviour.
\end{itemize}
\end{KeyConcept}

\begin{QuickRevision}
\begin{itemize}
\item Attitude $\rightarrow$ thought: forms attitude, then filters thoughts $\rightarrow$ confirmation bias (``we see things as we are'').
\item Attitude $\rightarrow$ behaviour: only a tendency; predicts behaviour better when it is STRONG, SPECIFIC and ACCESSIBLE, moderated by EXTERNAL FACTORS and PERCEIVED BEHAVIOURAL CONTROL.
\end{itemize}
\end{QuickRevision}

\section{Cognitive Dissonance}

When there is a large difference between what you are actually doing (behaviour) and what your attitude is, there is a sense of guilt and mental discomfort. This mental discomfort arising from misalignment between a strongly-held attitude and actual behaviour is called cognitive dissonance.

For example, we may have a strong attitude that we should study hard, but our actual behaviour is not that of studying --- so we should feel guilt/dissonance. If this misalignment is large and the attitude is strongly felt, we experience mental discomfort.

\begin{KeyConcept}
Cognitive dissonance = mental discomfort/guilt from a large misalignment between a strongly-held attitude and actual behaviour.
\end{KeyConcept}

\paragraph{Ways to resolve cognitive dissonance}
\begin{itemize}
\item Change the cognition: Alter what you believe. E.g., not studying despite a strong study-attitude, you add a new cognition --- ``so-and-so also didn't study much and still managed; I'll be fine'' --- to reduce the guilt.
\item Add a new cognition: Introduce an additional belief to reduce discomfort. E.g., ``I'll eat one chocolate despite the diet, but I'll run an extra kilometre'' --- adding a cognition so the attitude--behaviour dissonance is neutralised.
\item Change the behaviour: Bring behaviour back in line with the attitude (e.g., actually start studying).
\end{itemize}

\begin{ExampleBox}
\begin{itemize}
\item Diet vs.\ chocolate: ``I'll eat one chocolate but run an extra kilometre'' --- adding a cognition to remove dissonance.
\item Study attitude vs.\ no studying: rationalising (``others managed too'') to reduce guilt.
\end{itemize}
\end{ExampleBox}

\begin{CriticalPoint}
Resolving dissonance is important to avoid continuous mental discomfort --- done by changing cognition, adding a new cognition, or changing behaviour.
\end{CriticalPoint}

\section{Formation of Attitude}

Attitudes are formed with respect to an attitudinal object, through several mechanisms:

\subsection{Classical Conditioning (learning through association)}

We build attitudes by associating something with an experience over time. For example, songs or movies heard/watched while entering a relationship get associated with that positive experience; you later remember those songs/movies, and your attitude towards them stays positive because they are associated with a positive experience. Learning through association is classical conditioning.

\begin{KeyConcept}
Classical conditioning = learning through association (e.g., songs/movies tied to a positive experience gain a positive attitude).
\end{KeyConcept}

\subsection{Operant Conditioning (reward and punishment)}

Attitudes are inculcated through reward and punishment. A child who works hard and shows discipline is rewarded and encouraged by parents (praise, chocolate); a child who does something wrong is scolded (punishment). A child who steals and is punished forms a negative attitude towards stealing. Similarly, a child who touches fire and gets burnt has a negative experience/association and thereafter fears and avoids fire --- an attitude formed from experience/association.

\begin{KeyConcept}
Operant conditioning = attitude shaped by reward and punishment (e.g., stealing $\rightarrow$ punishment $\rightarrow$ negative attitude; touching fire $\rightarrow$ burn $\rightarrow$ avoidance).
\end{KeyConcept}

\subsection{Observational Learning, Culture, Media, Family}

We also learn much by observing others (observational learning), as discussed under value inculcation --- patriarchy and similar attitudes are learned this way. Culture plays a big role: we are born into a culture (e.g., the Bishnoi community's culture) and absorb its attitudes; a patriarchal society tends to produce patriarchal attitudes. Media plays a very important role, and family --- the family we are born into --- also contributes to attitude formation, as discussed under the role of value inculcation.

\begin{KeyConcept}
Other sources: observational learning; culture (e.g., Bishnoi community); media; family --- all contribute to attitude formation (as in value inculcation).
\end{KeyConcept}

\subsection{Values vs.\ Attitudes}

Values are our likes and preferences and are very broad in nature (e.g., sustainability, justice). Attitudes are specific --- towards a particular object. A broad value guides many specific attitudes: the value of sustainability guides a negative attitude towards plastic, a positive attitude towards trees, and a positive attitude towards water conservation. ``Justice'' is broad; ``justice for tribals'' or ``social justice for one section'' is more specific.

\begin{center}
\begin{tabularx}{\textwidth}{Y Y}
\rowcolor{AccentDark}
\thead{Values} & \thead{Attitudes} \\
\toprule
Broad in nature; our likes and preferences (e.g., sustainability, justice). & Specific; directed at a particular attitudinal object. \\
The broadest guiding layer. & The specific expression guided by values (e.g., sustainability $\rightarrow$ anti-plastic, pro-tree, pro-water-conservation). \\
\bottomrule
\end{tabularx}
\end{center}

\subsection{Moral, Immoral and Amoral Attitudes}

Just as values are of three types (moral, amoral, immoral), attitudes are also of three types:

\begin{itemize}
\item Amoral attitude: has no sense of right or wrong --- e.g., a positive attitude towards the colour black, or liking an Arijit Singh / Lata Mangeshkar song. No one would call these right or wrong.
\item Immoral attitude: e.g., greediness --- always thinking about money --- is an immoral attitude. Saying ``black people are bad people'' is a racist (immoral) attitude, unlike merely liking the colour black.
\item Moral attitude: what good values lead to. Cultivating good values ultimately leads to moral attitudes; otherwise it leads to immoral attitudes.
\end{itemize}

\begin{KeyConcept}
\begin{itemize}
\item Attitudes are of three types: moral, immoral, amoral (mirroring values).
\item Amoral = no right/wrong sense (liking black colour, a song); immoral = e.g., greed, racism; moral = product of good values.
\end{itemize}
\end{KeyConcept}

\begin{QuickRevision}
\begin{itemize}
\item Attitude forms via: classical conditioning (association), operant conditioning (reward/punishment), observational learning, culture, media, family.
\item Values (broad) guide attitudes (specific).
\item Attitudes may be moral, immoral or amoral.
\end{itemize}
\end{QuickRevision}
